\hypertarget{appendix-965d9260eeab}{}
\section*{Quaternion algebra and Hamilton product}
\section{Quaternion algebra and Hamilton product}\label{quaternion-algebra-and-hamilton-product}

\textbf{Curriculum depth:} Introduced · \textbf{First seen:} Phase7 · \textbf{Latest development:} Phase7

\subsection{Definition and mechanism}\label{definition-and-mechanism}

A quaternion has scalar w and vector u. Hamilton multiplication combines scalar product w\_a w\_b minus dot(u\_a,u\_b), and vector w\_a u\_b plus w\_b u\_a plus cross(u\_a,u\_b). It is associative and generally noncommutative; component order is a separate interface convention.

\textbf{Reconstruction cue:} i times j equals k, while j times i equals minus k.

\subsection{Prerequisites}\label{prerequisites}

\begin{itemize}
\tightlist
\item
  \href{}{Dot products and orthogonality (02-Foundations-II, appendix-8bc4e8fc547c)}
\item
  \href{}{Cross products and oriented area (02-Foundations-II, appendix-2391db8052b1)}
\end{itemize}

\subsection{Used by}\label{used-by}

\begin{itemize}
\tightlist
\item
  \hyperref[appendix-8e61ab68f6b2]{Quaternion rotation and double cover}
\item
  \hyperref[appendix-d1ea57347893]{Dual quaternions and rigid motion}
\end{itemize}

\subsection{Recall and next depth}\label{recall-and-next-depth}

Reconstruct the example, state its conventions and find an input that violates its preconditions. Deeper algorithms and measured production performance remain separate from the present model.

\subsection{Source and navigation}\label{source-and-navigation}

\href{Continuous\%20Notes/Computing\%20Foundations\%20-\%20From\%20Bits\%20to\%20C++\%20and\%20Systems\%20APIs\#Keep\%20the\%20promises\%20visible\%20while\%20moving\%20between\%20representations}{Continuous explanation} · \hyperref[appendix-fd090d982a15]{Coverage and source gap} · \hyperref[appendix-a103e7846d21]{Sources and verification}
